\documentclass[reqno]{amsart}
\usepackage[T1]{fontenc}
\usepackage{amsmath,amssymb,amsthm}
\usepackage{xr-hyper}
\usepackage[hidelinks]{hyperref}
\externaldocument[ms-]{20261125T002000--definite-theta-lifts__workbook}
\externaldocument[hp-]{20261125T020300--heights-periods-and-the-auxiliary-prime__workbook}
\newcounter{chap}
% Chapter numbers, generated from 20261125T012000--grh-workbook__guide.org
% by scripts/make-chapters.py.  Do not edit: change the order of the
% index and rerun the script (./preview-site.sh and CI do this).
%
% Usage in a sheet, after \newcounter{chap}:
%   % Chapter numbers, generated from 20261125T012000--grh-workbook__guide.org
% by scripts/make-chapters.py.  Do not edit: change the order of the
% index and rerun the script (./preview-site.sh and CI do this).
%
% Usage in a sheet, after \newcounter{chap}:
%   % Chapter numbers, generated from 20261125T012000--grh-workbook__guide.org
% by scripts/make-chapters.py.  Do not edit: change the order of the
% index and rerun the script (./preview-site.sh and CI do this).
%
% Usage in a sheet, after \newcounter{chap}:
%   \input{chapters}\chapnumber{<sheet ID>}
\makeatletter
\providecommand\chapters@set[2]{\@namedef{chapters@#1}{#2}}
\chapters@set{20261125T001000}{1}% zeros-to-ramp-sums
\chapters@set{20261125T001100}{2}% inert-primes-and-estimate-b
\chapters@set{20261125T005000}{3}% galois-fourier-and-averaging
\chapters@set{20261125T005500}{4}% columns-from-a-sphere
\chapters@set{20261125T000000}{5}% slope-method
\chapters@set{20261126T030100}{6}% a-twist-for-each-prime
\chapters@set{20261125T003000}{7}% the-chain-and-the-square-root
\chapters@set{20261125T010000}{8}% rank-lemma-and-zero-estimate
\chapters@set{20261125T010500}{9}% poincare-series-and-near-orthogonality
\chapters@set{20261124T223000}{10}% elliptic-curves-for-the-proof
\chapters@set{20261125T009500}{11}% fock-model-and-theta-intertwiners
\chapters@set{20261124T223100}{12}% complex-multiplication-and-inert-reduction
\chapters@set{20261125T008000}{13}% one-cm-determinant
\chapters@set{20261124T214500}{14}% supersingular-theta-divisibility
\chapters@set{20261124T214600}{15}% the-isotropic-plane
\chapters@set{20261125T013700}{16}% pairings-on-a-formal-group-of-height-two
\chapters@set{20261125T013500}{17}% hecke-nondegeneracy-on-the-tree
\chapters@set{20261126T030200}{18}% the-bound-is-attained
\chapters@set{20261125T009000}{19}% kronecker-limit-and-chowla-selberg
\chapters@set{20261125T002000}{20}% definite-theta-lifts
\chapters@set{20261125T015500}{21}% siegel-weil-in-the-smallest-case
\chapters@set{20261125T002100}{22}% definite-theta-lifts-and-norms
\chapters@set{20261125T016500}{23}% hecke-operators-and-theta-lifts
\chapters@set{20261125T016700}{24}% divisibility-along-the-chain
\chapters@set{20261125T004000}{25}% denominators-and-sha
\chapters@set{20261125T004100}{26}% tate-twists-and-sha
\chapters@set{20261126T020430}{27}% the-minus-part-by-gauss-sums
\chapters@set{20261125T014100}{28}% congruence-modules
\chapters@set{20261125T014500}{29}% eisenstein-congruences-and-selmer-classes
\chapters@set{20261125T014700}{30}% pseudocharacters-and-the-length-bound
\chapters@set{20261125T011000}{31}% galois-representations-black-box
\chapters@set{20261126T030300}{32}% the-data
\chapters@set{20261126T020210}{33}% bessel-poincare-series-and-their-gram-operators
\chapters@set{20261126T030400}{34}% the-seesaw-identity
\chapters@set{20261126T030500}{35}% divisibility-of-the-composite-lift
\chapters@set{20261126T030600}{36}% theorem-a-short-forms
\chapters@set{20261126T030700}{37}% the-determinant-of-a-chain
\chapters@set{20261126T030800}{38}% the-cost-of-one-lift
\chapters@set{20261126T030900}{39}% denominators-of-projections-to-lifts
\chapters@set{20261126T031000}{40}% twists-with-few-congruences
\chapters@set{20261126T031100}{41}% theorem-c-and-the-main-inequality
\providecommand\chapnumber[1]{%
  \@ifundefined{chapters@#1}%
    {\PackageError{chapters}{Sheet #1 is not listed in the index}%
      {List it in 20261125T012000--grh-workbook__guide.org\MessageBreak
       and run scripts/make-chapters.py.}}%
    {\edef\thechap{\@nameuse{chapters@#1}}}}
\makeatother
\chapnumber{<sheet ID>}
\makeatletter
\providecommand\chapters@set[2]{\@namedef{chapters@#1}{#2}}
\chapters@set{20261125T001000}{1}% zeros-to-ramp-sums
\chapters@set{20261125T001100}{2}% inert-primes-and-estimate-b
\chapters@set{20261125T005000}{3}% galois-fourier-and-averaging
\chapters@set{20261125T005500}{4}% columns-from-a-sphere
\chapters@set{20261125T000000}{5}% slope-method
\chapters@set{20261126T030100}{6}% a-twist-for-each-prime
\chapters@set{20261125T003000}{7}% the-chain-and-the-square-root
\chapters@set{20261125T010000}{8}% rank-lemma-and-zero-estimate
\chapters@set{20261125T010500}{9}% poincare-series-and-near-orthogonality
\chapters@set{20261124T223000}{10}% elliptic-curves-for-the-proof
\chapters@set{20261125T009500}{11}% fock-model-and-theta-intertwiners
\chapters@set{20261124T223100}{12}% complex-multiplication-and-inert-reduction
\chapters@set{20261125T008000}{13}% one-cm-determinant
\chapters@set{20261124T214500}{14}% supersingular-theta-divisibility
\chapters@set{20261124T214600}{15}% the-isotropic-plane
\chapters@set{20261125T013700}{16}% pairings-on-a-formal-group-of-height-two
\chapters@set{20261125T013500}{17}% hecke-nondegeneracy-on-the-tree
\chapters@set{20261126T030200}{18}% the-bound-is-attained
\chapters@set{20261125T009000}{19}% kronecker-limit-and-chowla-selberg
\chapters@set{20261125T002000}{20}% definite-theta-lifts
\chapters@set{20261125T015500}{21}% siegel-weil-in-the-smallest-case
\chapters@set{20261125T002100}{22}% definite-theta-lifts-and-norms
\chapters@set{20261125T016500}{23}% hecke-operators-and-theta-lifts
\chapters@set{20261125T016700}{24}% divisibility-along-the-chain
\chapters@set{20261125T004000}{25}% denominators-and-sha
\chapters@set{20261125T004100}{26}% tate-twists-and-sha
\chapters@set{20261126T020430}{27}% the-minus-part-by-gauss-sums
\chapters@set{20261125T014100}{28}% congruence-modules
\chapters@set{20261125T014500}{29}% eisenstein-congruences-and-selmer-classes
\chapters@set{20261125T014700}{30}% pseudocharacters-and-the-length-bound
\chapters@set{20261125T011000}{31}% galois-representations-black-box
\chapters@set{20261126T030300}{32}% the-data
\chapters@set{20261126T020210}{33}% bessel-poincare-series-and-their-gram-operators
\chapters@set{20261126T030400}{34}% the-seesaw-identity
\chapters@set{20261126T030500}{35}% divisibility-of-the-composite-lift
\chapters@set{20261126T030600}{36}% theorem-a-short-forms
\chapters@set{20261126T030700}{37}% the-determinant-of-a-chain
\chapters@set{20261126T030800}{38}% the-cost-of-one-lift
\chapters@set{20261126T030900}{39}% denominators-of-projections-to-lifts
\chapters@set{20261126T031000}{40}% twists-with-few-congruences
\chapters@set{20261126T031100}{41}% theorem-c-and-the-main-inequality
\providecommand\chapnumber[1]{%
  \@ifundefined{chapters@#1}%
    {\PackageError{chapters}{Sheet #1 is not listed in the index}%
      {List it in 20261125T012000--grh-workbook__guide.org\MessageBreak
       and run scripts/make-chapters.py.}}%
    {\edef\thechap{\@nameuse{chapters@#1}}}}
\makeatother
\chapnumber{<sheet ID>}
\makeatletter
\providecommand\chapters@set[2]{\@namedef{chapters@#1}{#2}}
\chapters@set{20261125T001000}{1}% zeros-to-ramp-sums
\chapters@set{20261125T001100}{2}% inert-primes-and-estimate-b
\chapters@set{20261125T005000}{3}% galois-fourier-and-averaging
\chapters@set{20261125T005500}{4}% columns-from-a-sphere
\chapters@set{20261125T000000}{5}% slope-method
\chapters@set{20261126T030100}{6}% a-twist-for-each-prime
\chapters@set{20261125T003000}{7}% the-chain-and-the-square-root
\chapters@set{20261125T010000}{8}% rank-lemma-and-zero-estimate
\chapters@set{20261125T010500}{9}% poincare-series-and-near-orthogonality
\chapters@set{20261124T223000}{10}% elliptic-curves-for-the-proof
\chapters@set{20261125T009500}{11}% fock-model-and-theta-intertwiners
\chapters@set{20261124T223100}{12}% complex-multiplication-and-inert-reduction
\chapters@set{20261125T008000}{13}% one-cm-determinant
\chapters@set{20261124T214500}{14}% supersingular-theta-divisibility
\chapters@set{20261124T214600}{15}% the-isotropic-plane
\chapters@set{20261125T013700}{16}% pairings-on-a-formal-group-of-height-two
\chapters@set{20261125T013500}{17}% hecke-nondegeneracy-on-the-tree
\chapters@set{20261126T030200}{18}% the-bound-is-attained
\chapters@set{20261125T009000}{19}% kronecker-limit-and-chowla-selberg
\chapters@set{20261125T002000}{20}% definite-theta-lifts
\chapters@set{20261125T015500}{21}% siegel-weil-in-the-smallest-case
\chapters@set{20261125T002100}{22}% definite-theta-lifts-and-norms
\chapters@set{20261125T016500}{23}% hecke-operators-and-theta-lifts
\chapters@set{20261125T016700}{24}% divisibility-along-the-chain
\chapters@set{20261125T004000}{25}% denominators-and-sha
\chapters@set{20261125T004100}{26}% tate-twists-and-sha
\chapters@set{20261126T020430}{27}% the-minus-part-by-gauss-sums
\chapters@set{20261125T014100}{28}% congruence-modules
\chapters@set{20261125T014500}{29}% eisenstein-congruences-and-selmer-classes
\chapters@set{20261125T014700}{30}% pseudocharacters-and-the-length-bound
\chapters@set{20261125T011000}{31}% galois-representations-black-box
\chapters@set{20261126T030300}{32}% the-data
\chapters@set{20261126T020210}{33}% bessel-poincare-series-and-their-gram-operators
\chapters@set{20261126T030400}{34}% the-seesaw-identity
\chapters@set{20261126T030500}{35}% divisibility-of-the-composite-lift
\chapters@set{20261126T030600}{36}% theorem-a-short-forms
\chapters@set{20261126T030700}{37}% the-determinant-of-a-chain
\chapters@set{20261126T030800}{38}% the-cost-of-one-lift
\chapters@set{20261126T030900}{39}% denominators-of-projections-to-lifts
\chapters@set{20261126T031000}{40}% twists-with-few-congruences
\chapters@set{20261126T031100}{41}% theorem-c-and-the-main-inequality
\providecommand\chapnumber[1]{%
  \@ifundefined{chapters@#1}%
    {\PackageError{chapters}{Sheet #1 is not listed in the index}%
      {List it in 20261125T012000--grh-workbook__guide.org\MessageBreak
       and run scripts/make-chapters.py.}}%
    {\edef\thechap{\@nameuse{chapters@#1}}}}
\makeatother

\makeatletter
\@ifundefined{chapters@20261126T020310}{\renewcommand{\thechap}{?}}{\chapnumber{20261126T020310}}
\makeatother
\renewcommand{\thesection}{\thechap.\arabic{section}}
\theoremstyle{definition}
\newtheorem{problem}{Problem}[section]
\theoremstyle{plain}
\newtheorem{theorem}{Theorem}[section]
\newtheorem{lemma}[theorem]{Lemma}
\theoremstyle{remark}
\newtheorem*{leanremark}{In the Lean formalization}
\newcommand{\mass}{\operatorname{mass}}
\newcommand{\rd}{\operatorname{rd}}
\newcommand{\GL}{\mathrm{GL}}
\newcommand{\Q}{\mathbb{Q}}
\newcommand{\F}{\mathbb{F}}
\newcommand{\Tc}{T^{\mathrm{cr}}}
\newcommand{\zr}{z^{\mathrm{row}}}
\newcommand{\zc}{z^{\mathrm{col}}}
\setlength{\emergencystretch}{2em}
\title{Counting forms and the critical level}
\author{}
\date{November 26, 2026}
\begin{document}
\maketitle
\refstepcounter{chap}\label{chap}

\paragraph{Prerequisites.}
The mass of a genus, Weyl's dimension formula and the count identity of
Chapter~\ref{ms-chap}; the dimension budget, the level indices and the
interval for the auxiliary prime in \S\ref{hp-sec:prime}.  We use the
order of $\GL_n(\F_q)$, the Euler product of the Dedekind zeta function
of a number field, and $\log(1+t)\le t$.

This chapter develops the input (A6) of Chapter~\ref{hp-chap} from the
Lean formalization \cite{Lean}, where the published expositions stop.
There the dimension budget is not an inequality: the rank of the source
lattice and the budget are both given exactly by the mass formula, and
their ratio is a product over the embeddings of a \emph{critical level}
divided by the auxiliary prime.  The prime is then pinned to the
geometric mean of the critical levels, and the lower bound (A6) becomes a
lower bound for each critical level, which is a Vandermonde estimate.

\section{Two dimensions}\label{sec:two}

How many forms does each side of the final pairing carry?  Let $F$ be a
totally real number field of degree $f$, and $E\supset F$ a CM field
unramified over $F$ at every finite prime, so that $|d_E|=d_F^2$ for the
discriminants.  Let $\chi$ be the quadratic character of $F$ attached
to $E$, write $\rd_E=|d_E|^{1/(2f)}$ and $v=\rd_E/2$ as in
Chapter~\ref{hp-chap}, and let $\Phi$ be a CM type, a set of $f$ complex
embeddings $\rho$ of $E$.  For $M\ge1$ let $\mass(M)$ be the mass of the
genus of the standard lattice $\mathcal O_E^M$ with
$h(x,y)=\sum_i\bar x_iy_i$, the sum of $1/|\mathrm{Aut}|$ over its classes.

\begin{theorem}[Mass ratio]\label{in:mass}
For $M\ge3$,
\[
 \frac{\mass(M)}{\mass(M-1)}
 =\bigl((M-1)!\,(v/\pi)^M\bigr)^f(2v)^{-f/2}\,L(M),
 \qquad L(M)=\prod_{\mathfrak P}e_{\mathfrak P}(M),
\]
the product over the primes $\mathfrak P$ of $E$, with
$e_{\mathfrak P}(M)=(1-(-1)^M N\mathfrak P^{-M/2})^{-1}$ if
$\bar{\mathfrak P}=\mathfrak P$ and
$e_{\mathfrak P}(M)=(1-N\mathfrak P^{-M})^{-1/2}$ otherwise.
\end{theorem}
\noindent\emph{Status}: this follows from the exact mass formula of Gan, Hanke and
Yu \cite{GHY} and the functional equation.  In \cite{Lean} it is proved
from a statement on representations of one by hermitian lattices; we have
not followed that proof.
Simplest incarnation: the Gaussian ratio
$M_r/M_{r-1}=(r-1)!\,\pi^{-r}\kappa_r$ of \eqref{ms-eq:kappa}, where
$\mathbb Q(\mathrm i)$ is ramified at two and so lies outside the present
hypotheses.

\begin{problem}\label{pr:euler}
For $M\ge2$, $L(M)=\zeta_F(M)$ for even $M$, and $L(M)=L_F(M,\chi)$ for odd $M$.
\end{problem}

\begin{problem}\label{pr:Lbound}
For $M\ge2$, $|\log L(M)|\le\log\zeta_F(M)\le f\log\zeta(M)$.
\end{problem}

Fix even $m=2s+2\ge2$, $k\ge m+3$ and $r=k+m-1$, and two primes
$g_0\ne g_1$ split completely in $E$.  The level indices of
\S\ref{hp-sec:prime} are $I_\Pi=(I_\Pi(g_0)I_\Pi(g_1))^f$ and
$I_\Sigma=(I_\Sigma(g_0)I_\Sigma(g_1))^f$ with
$I_\Pi(q)=\prod_{i=1}^k(q^i-1)$ and
$I_\Sigma(q)=q^k\prod_{i=1}^r(q^i-1)$.

\begin{problem}\label{pr:glcard}
$|\GL_n(\F_q)|=q^{n(n-1)/2}\prod_{i=1}^n(q^i-1)$.  With
$N=s(2k+2s+1)$ one has $k^2+N+km+m(m-1)/2=r^2$, and hence
$I_\Sigma(q)=|\GL_r(\F_q)|/q^{N+k(k-1)/2}$.
\end{problem}

\begin{problem}\label{pr:index}
$I_\Sigma(q)/I_\Pi(q)=q^{km+m(m-1)/2}\prod_{j=k+1}^{r}(1-q^{-j})$.
For $k=5$ and $m=2$ this is $q^5(q^6-1)$.
\end{problem}

Recall from \eqref{ms-eq:weyl} the Vandermonde product
$V_n(z)=\prod_{i<i'}(z_i-z_{i'})$ and $D_n(z)=V_n(z)/\prod_{i=0}^{n-1}i!$.
For each $\rho\in\Phi$ let $\zr_\rho\in\mathbb R^r$ be the row parameter
and, for each word $u$, $\zc_{u\rho}\in\mathbb R^k$ the column parameter,
both strictly decreasing.  Let $T_u$ be the auxiliary prime of the word,
let $d_u$ be the rank of the source lattice $N_u$, and $D^*$ the
dimension budget.

\begin{theorem}[The two counts]\label{in:counts}
\[
 d_u=\mass(k)\,I_\Pi\prod_{\rho\in\Phi}D_k(\zc_{u\rho})\,(T_u^k-1)^f,
 \qquad
 D^*=\mass(r)\,I_\Sigma\prod_{\rho\in\Phi}D_r(\zr_\rho).
\]
\end{theorem}
\noindent\emph{Status}: the first is the class number formula for the column
group at level $T_u$ together with Weyl's dimension formula; the second is
the volume of the row group times Harish-Chandra's formal degrees
\cite{HC}, which are given by Weyl's formula.  Both are proved in
\cite{Lean} from these inputs; we have read the statements, not the
proofs.  In \cite[Hyp.~15.23]{Adr} the budget is only bounded below, with
a factor $(2/3)^f$.

For each word $u$ and $\rho\in\Phi$ define the \emph{critical level}
$\Tc_{u\rho}>0$ by
\[
 (\Tc_{u\rho})^k=(v/\pi)^{(m-1)k+m(m-1)/2}\,\rd_E^{-(m-1)/2}\,
 (g_0g_1)^{km+m(m-1)/2}\,\frac{V_r(\zr_\rho)}{V_k(\zc_{u\rho})},
\]
and put
\[
 E_a=\prod_{M=k+1}^{r}L(M)\cdot
 \prod_{i=0,1}\prod_{j=k+1}^{r}(1-g_i^{-j})^f\cdot(1-T_u^{-k})^{-f}.
\]

\begin{problem}\label{pr:telescope}
$\sum_{M=k+1}^{r}M=(m-1)k+m(m-1)/2$, and
\[
 \frac{\mass(r)}{\mass(k)}=\Bigl((v/\pi)^{(m-1)k+m(m-1)/2}\,
 \rd_E^{-(m-1)/2}\prod_{i=k}^{r-1}i!\Bigr)^f\prod_{M=k+1}^{r}L(M).
\]
\end{problem}

\begin{problem}\label{pr:factorials}
$\prod_{i=k}^{r-1}i!\cdot D_r(z)/D_k(z')=V_r(z)/V_k(z')$ for all
$z\in\mathbb R^r$ and $z'\in\mathbb R^k$ with $V_k(z')\ne0$.
\end{problem}

Combining Theorem~\ref{in:counts} with Problems~\ref{pr:index},
\ref{pr:telescope} and~\ref{pr:factorials} gives the following.

\begin{theorem}[The count ratio]\label{thm:count}
For every word $u$,
\[
 \frac{D^*}{d_u}=E_a\prod_{\rho\in\Phi}\Bigl(\frac{\Tc_{u\rho}}{T_u}\Bigr)^k .
\]
\end{theorem}

No factorial survives, as in the Gaussian count identity
(Theorem~\ref{ms-thm:count}).  What remains of the masses is the power of
$v/\pi$ and the $L$-values in $E_a$.

\section{The prime and the critical levels}\label{sec:prime}

How close is $T_u$ to the critical levels?  The factor $E_a$ is
exponentially close to one in $k$.

\begin{problem}\label{pr:zeta-tail}
For $M\ge3$, $\log\zeta(M)\le2^{1-M}$.
\end{problem}

\begin{problem}\label{pr:Ea}
If $g_0,g_1\ge4$ and $T_u\ge2$, then $|\log E_a|\le8f/2^k$.
\end{problem}

The prime $T_u$ is chosen so that $d_u$ matches the budget up to the
frame factor $\Gamma_0\ge1$ of \S\ref{hp-sec:prime}, the same for all
words: for a constant $C_0\ge1$,
\begin{equation}\label{eq:window}
 C_0^{-1}\Gamma_0^{-2}D^*\le d_u\le C_0\Gamma_0^{-2}D^*.
\end{equation}
The formalization takes $C_0=3$.  The prime itself is part of its data,
subject to the interval $Z_u\le T_u\le(1+1/(fk))Z_u$ of
Problem~\ref{hp-pr:interval}.

\begin{problem}\label{pr:window}
Assume \eqref{eq:window}, $8f/2^k\le\log2$ and $\log\Gamma_0\le k/8$,
together with the hypotheses of Problem~\ref{pr:Ea}.  Then
\[
 (2C_0)^{-1/k}\le\prod_{\rho\in\Phi}\frac{\Tc_{u\rho}}{T_u}
 \le(2C_0)^{1/k}e^{1/4}.
\]
\end{problem}

In the application, $\Gamma_0=\exp(C_\Gamma m^5fy)$ with $y=x^\varepsilon$, and
the size condition $8C_\Gamma m^5fy\le k$ gives $\log\Gamma_0\le k/8$.
Thus $\log T_u$ is the mean of $\log\Tc_{u\rho}$ over $\Phi$ up to
$O(1/f)$.  The same comparison bounds the ratio of the primes of two
words.

\begin{problem}\label{pr:ratio}
Let $u,u'$ be words satisfying \eqref{eq:window}, with
$|\log E_a|\le8f/2^k$ for both and $k\ge5$.  If
$\Tc_{u\rho}\le C_1\Tc_{u'\rho}$ for every $\rho$, then
$T_u\le eC_0C_1T_{u'}$.
\end{problem}

\section{The critical level from below}\label{sec:below}

Which part of $\Tc_{u\rho}$ is large?  Fix $\rho$ and drop it from the
notation.  The row parameter is
$\zr=(\alpha_1,\dots,\alpha_{k-1},-b_1,\dots,-b_m)$ with
\[
 \alpha_i=\alpha_{k-1}+4\Delta(k-1-i),\qquad
 \alpha_{k-1}\ge0,\quad \alpha_1\le4k\Delta,\quad\Delta\ge1,
\]
and nodes $1\le b_1$, $b_{j+1}\ge b_j+1$.  A word is
$\delta\in\{0,1\}^{k-1}$ with $\delta_{k-1}=0$, and the column parameter
is $\zc=(\gamma_1,\dots,\gamma_k)$ with
\[
 \gamma_i=\alpha_{i-1}-\tfrac12-\delta_{i-1}\Delta\quad(2\le i\le k),
 \qquad \gamma_1-\alpha_1\ge4k\Delta,
\]
and we put $\eta=\gamma_1-\alpha_1$.
This is the case of a place $\rho$ that is not a top.  At a top the
last entry is replaced by $\gamma_k=-b'$ with $b_1/2\le b'\le b_1-3$, and
we put $\eta=\gamma_1b'/(k\Delta)$ instead.

\begin{problem}\label{pr:interlace}
$\gamma_1>\alpha_1>\gamma_2>\alpha_2>\dots>\alpha_{k-1}>\gamma_k$, and
$\zr$ is strictly decreasing.
\end{problem}

\begin{problem}\label{pr:bracket}
Let $B>0$ with
$B^k=(v/\pi)^{m(m-1)/2}\rd_E^{-(m-1)/2}(g_0g_1)^{m(m-1)/2}$.  Then
\[
 \Tc=(v/\pi)^{m-1}(g_0g_1)^m\,B\,
 \bigl(V_r(\zr)/V_k(\zc)\bigr)^{1/k},
\]
and if $\rd_E\ge1$ and $g_0g_1\ge2\pi$, then
$1\le B\le\exp\bigl(m^2\log(g_0g_1\rd_E)/(2k)\bigr)$.
\end{problem}

\begin{problem}\label{pr:row}
\[
 V_r(\zr)=V_{k-1}(\alpha)\prod_{i,j}(\alpha_i+b_j)\prod_{j<j'}(b_{j'}-b_j)
 \ge V_{k-1}(\alpha)\,(b_1\cdots b_m)^{k-1}.
\]
\end{problem}

\begin{problem}\label{pr:split}
$V_k(\zc)=\prod_{i=2}^k(\gamma_1-\gamma_i)\cdot
V_{k-1}(\gamma_2,\dots,\gamma_k)$, and
\[
 V_{k-1}(\gamma_2,\dots,\gamma_k)=R_\delta\,V_{k-1}(\alpha),\qquad
 R_\delta=\prod_{1\le i<i'\le k-1}
 \Bigl(1+\frac{\delta_{i'}-\delta_i}{4(i'-i)}\Bigr).
\]
\end{problem}

\begin{problem}\label{pr:R}
$R_\delta\le e^{(k-2)/4}$.
\end{problem}

The formalization also proves $R_\delta\ge e^{-k/4}$; we do not need it.

\begin{problem}\label{pr:edge}
At a non-top place, $\eta\le\gamma_1-\gamma_i\le2\eta$ for $2\le i\le k$.  Hence
$V_k(\zc)\le(2\eta)^{k-1}e^{k/4}V_{k-1}(\alpha)$.
\end{problem}

\begin{problem}\label{pr:constant}
$2\cdot0.38\cdot e^{1/4}<1$.  Let $\beta=b_1\cdots b_m$, let $b\ge1$ with
$b_j\le b^2$ for all $j$, and put $e_0=\exp(-2m\log b/k)$.  Then
$e_0^k\beta\le1$, and
$(0.38\,e_0\beta/\eta)^k\,V_k(\zc)\le\tfrac12V_{k-1}(\alpha)\beta^{k-1}$.
\end{problem}

By Problems~\ref{pr:bracket}, \ref{pr:row} and~\ref{pr:constant}, at a
non-top place $(V_r(\zr)/V_k(\zc))^{1/k}\ge0.38\,e_0\beta/\eta$, and
$B\ge1$.  The top case runs the same way with a separate Vandermonde
count for the last column entry.

\begin{theorem}[Lower bound for the critical level]\label{thm:below}
Let $k\ge4$, $b\ge1$, $b_j\le b^2$ for all $j$, $\rd_E\ge1$ and
$g_0g_1\ge2\pi$.  At a non-top place
\[
 \Tc\ge c_1\,(g_0g_1)^m(v/\pi)^{m-1}\,\frac{b_1\cdots b_m}{\eta},
 \qquad c_1=0.38\exp(-2m\log b/k).
\]
At a top, the same holds with $c_1$ divided by $1+4k\Delta/b'$, under
$40mk\Delta\le b_1$ and $20m\log b\le k$.
\emph{Status}: the non-top case is Problems~\ref{pr:bracket}--\ref{pr:constant};
the top case is proved in \cite{Lean}, and we have not checked its
Vandermonde count by hand.
\end{theorem}

\begin{problem}\label{pr:c1}
If $40mk\Delta\le b_1$ and $20m\log b\le k$, then $c_1\ge1/5$ in both
cases of Theorem~\ref{thm:below}.
\end{problem}

The word enters $\Tc$ only through $V_k(\zc)$.  The formalization also
proves that, for $40mk\Delta\le b_1$ and any two words $\delta,\delta'$,
$|\log(V_k(\zc_{\delta'})/V_k(\zc_\delta))|\le k/2+1$ \cite{Lean}.

\begin{problem}\label{pr:words}
Under that bound and $k\ge4$, the critical levels of two words differ by
a factor between $e^{-3/4}$ and $e^{3/4}$.
\end{problem}

With Problem~\ref{pr:ratio} this gives $T_u\le e^{7/4}C_0T_{u'}$ for any
two words, comparing the critical levels at each place.

\section{The auxiliary prime}\label{sec:A6}

Combining Problem~\ref{pr:window} with Theorem~\ref{thm:below} and
Problem~\ref{pr:c1} gives the following.

\begin{theorem}[Lower bound for the prime]\label{thm:A6}
Under the hypotheses of Problems~\ref{pr:window} and~\ref{pr:c1} and of
Theorem~\ref{thm:below} at every $\rho\in\Phi$, with nodes $b_{\rho j}$
and $\eta_{u\rho}$ at $\rho$,
\begin{multline*}
 \log T_u\ge m\log(g_0g_1)+(m-1)\log\frac v\pi
 +\frac1f\sum_{\rho\in\Phi}\Bigl(\sum_{j=1}^m\log b_{\rho j}
 -\log\eta_{u\rho}\Bigr)\\
 -\log5-\frac1f\Bigl(\frac{\log(2C_0)}k+\frac14\Bigr).
\end{multline*}
\end{theorem}

\begin{problem}\label{pr:A6}
Theorem~\ref{thm:A6} follows from Problem~\ref{pr:window},
Theorem~\ref{thm:below} and Problem~\ref{pr:c1}.
\end{problem}

\begin{problem}\label{pr:compare}
Let $y=\log b$, $g=\log(g_0g_1)$, and suppose $\log(v/\pi)\le y$,
$b_{\rho j}\le b^2$ and $\eta_{u\rho}\le b_{\rho1}b^{C_u}$.  Then
Theorem~\ref{thm:A6} gives
\[
 \log T_u\ge\sum_{j=1}^m\frac1f\sum_{\rho\in\Phi}
 \log\frac{b_{\rho j}v}{\pi}+mg-(3+C_u)y-\log5
 -\frac1f\Bigl(\frac{\log(2C_0)}k+\frac14\Bigr).
\]
\end{problem}

This is Theorem~\ref{hp-in:A6}, the form of (A6) with the nodes
$b_{\rho j}$ in place of $w_{j\rho}$ and an error of order $y$.  In
\cite[\S15.5]{Adr} the nodes satisfy $b_{\rho j}<w_{j\rho}\le
b_{\rho j}(1+1/x)$, so the two forms differ by at most $m/x$.  The
formalization bounds the first column entry by $\gamma_1\le b^{C_u}$, which
gives $\eta_{u\rho}\le b_{\rho1}b^{C_u}$ at a top; elsewhere it imposes
$\eta_{u\rho}\le8k\Delta$, which is at most $b_{\rho1}$ when
$40mk\Delta\le b_{\rho1}$.

\begin{leanremark}\sloppy
The mass ratio is \texttt{SiegelMassRatioCM}, proved as
\texttt{massRatio\_all\_unramified\_cmFields}, with right side
\texttt{massRatioRightSide} and $L(M)$ \texttt{lValEta}; the bound of
Problem~\ref{pr:Lbound} is \texttt{log\_l\_value\_field\_bound}.
Theorem~\ref{in:counts} is \texttt{FamilyMidVolume}, proved by
\texttt{admissible\_ranks\_of\_mass\_alt} through \texttt{Dstar\_mass} and
\texttt{admissibleRanks\_of\_mass\_formula}, with formal degrees
\texttt{DsigHC}.  The critical level is \texttt{Tcrit}.
Theorem~\ref{thm:count} is \texttt{Sz.PropA}, used in
\texttt{degree\_ratio\_product\_formula}; Problems~\ref{pr:Ea},
\ref{pr:window} and~\ref{pr:ratio} are \texttt{Sz.EaSize},
\texttt{Sz.CorA} and \texttt{level\_ratio\_bound}.  Section~\ref{sec:below}
follows \texttt{BracketBound}, \texttt{VandLamLower}, \texttt{VandNonTop},
\texttt{VandTop}, \texttt{EdgeFactors} and \texttt{TcritLower} (proved as
\texttt{critical\_threshold\_lower\_bound}, used in
\texttt{critical\_level\_lower\_bound}), and the word comparison is
\texttt{TcritWords}.  The window \eqref{eq:window} is part of
\texttt{StandingNumericalData}; the interval for $T_u$ is
\texttt{levelPrime\_window}.  Downstream, the formalization uses
$T_u\ge(g_0g_1v/\pi)^{m-2}$ (\texttt{TFrameAt}) and
$b^{m/2}\le T_u\le b^{(4C_g+13)m}$ (\texttt{levelPrime\_size})
\cite{Lean}.
\end{leanremark}

\emph{Status}: we have not traced where the formalization uses a lower
bound for $\log T_u$ in the mean slope inequality, nor how it constructs a
prime in the window; we located only the consequences named above.
Theorem~\ref{thm:A6} is our combination of its lemmas.

\section*{Hints}
\subsection*{\S\ref{sec:two}}
\begin{description}
\item[\ref{pr:euler}] Since $E/F$ is unramified, $\bar{\mathfrak P}=\mathfrak P$
exactly when $\mathfrak P$ lies over an inert prime of $F$.  Group the
factors over each prime of $F$.
\item[\ref{pr:Lbound}] Compare each Euler factor with
$(1-N\mathfrak p^{-M})^{\pm1}$, and the primes of $F$ over $p$ with
$(1-p^{-M})^{-f}$.
\item[\ref{pr:glcard}] Count ordered bases.  For the identity, expand
$r^2$ with $r=k+2s+1$.
\item[\ref{pr:telescope}] Multiply Theorem~\ref{in:mass} over $M$; use $2v=\rd_E$.
\end{description}
\subsection*{\S\ref{sec:prime}}
\begin{description}
\item[\ref{pr:zeta-tail}] Bound $\sum_{n\ge3}n^{-M}$ by an integral.
\item[\ref{pr:Ea}] Use $-\log(1-t)\le2t$ for $0\le t\le\frac12$ and
Problems~\ref{pr:Lbound} and~\ref{pr:zeta-tail}.
\item[\ref{pr:window}] Theorem~\ref{thm:count} with
$\frac12\le E_a\le2$.
\item[\ref{pr:ratio}] Compare $T_u^{kf}\prod_\rho(\Tc_{u\rho}/T_u)^k$ for $u$ and $u'$;
\eqref{eq:window} gives $d_u\le C_0^2d_{u'}$.
\end{description}
\subsection*{\S\ref{sec:below}}
\begin{description}
\item[\ref{pr:bracket}] Substitute $\rd_E=2v$ in $B^k$.
\item[\ref{pr:R}] Group the pairs by $d=i'-i$.
\item[\ref{pr:edge}] $\gamma_k$ is the smallest column entry.  Then use
Problems~\ref{pr:split} and~\ref{pr:R}.
\item[\ref{pr:constant}] Insert Problem~\ref{pr:edge} and use $\eta\ge1$.
\item[\ref{pr:c1}] At a top, $b'\ge b_1/2$; and $e^{-t}\ge1-t$.
\end{description}
\subsection*{\S\ref{sec:A6}}
\begin{description}
\item[\ref{pr:A6}] Take logarithms in Problem~\ref{pr:window} and average
Theorem~\ref{thm:below} over $\Phi$.
\end{description}

\begin{thebibliography}{GHY}
\bibitem[Adr]{Adr}
J.~Adriazola, \emph{The Alp\"oge proof of the Riemann hypothesis}, manuscript, revised 24 November 2026, \url{https://github.com/jimjamzola/RH-Alpoge-Reconstruction}.
\bibitem[GHY]{GHY} W.~T.~Gan, J.~P.~Hanke and J.-K.~Yu, \emph{On an exact mass formula of Shimura}, Duke Math.\ J.\ \textbf{107} (2001), 103--133.
\bibitem[HC]{HC} Harish-Chandra, \emph{Discrete series for semisimple Lie groups.~II}, Acta Math.\ \textbf{116} (1966), 1--111.
\bibitem[Lean]{Lean} The Lean community, \emph{riemann-hypothesis}, a Lean formalization of Alp\"oge's argument, \url{https://github.com/leanprover-community/riemann-hypothesis}.
\end{thebibliography}
\end{document}
